% =====================================================================
\section{Critical Analysis of the Base Scheme}\label{sec:findings}
% =====================================================================
Each finding states \emph{what} fails, \emph{where} in the paper, \emph{why}, and the \emph{impact} on the paper's claims.
Where possible we confirmed the finding with a script; such cases are marked \emph{verified}.

% ---------------------------------------------------------------------
\subsection{Finding 1: the scheme is not infrastructure-less}
\begin{flawbox}[title={Authentication requires a live satellite link to the TA}]
\begin{finding}\label{f:online}
In Step GA1 the TUAV uploads every request to the TA. In Step GA2 the TA computes $\sigma_i,\psi_i,\mu_i,\delta_i$, which include all
pairings and need the secret $s_i$. The verification equation~\eqref{eq:ga3} of Step GA3 has no inputs other than these values. Nothing in the paper
gives a local fallback.
\end{finding}
\textbf{Why it matters.} The paper's own system model says the TA is reached through ``space networks'' (satellite), because the
terrestrial backhaul is destroyed. In the paper's scenario this link is shared, congested and possibly jammed. If $A$ is the probability that the
TA link is usable, the probability that an authentication round succeeds is exactly $A$. Its latency also includes one satellite round trip
(typically about 30--60\,ms for LEO and about 500--600\,ms for GEO).

\textbf{Impact.} When the link is lost, authentication does not slow down, it \emph{stops}. The dependency on infrastructure was not removed; it
was moved from a roadside pole to a satellite link that nobody in the field controls.
\end{flawbox}

% ---------------------------------------------------------------------
\subsection{Finding 2: the TA can impersonate every UAV}
Look again at Step US2, equation~\eqref{eq:gamma}. The last term uses $h_1(r_i)$, and Step US1 defined $\xi_i=h_1(r_i)$, which is sent in clear
in Step US3. Substituting gives
\begin{equation}
  \Gamma_i=\big[s_i\,h_3(ID_i,ts_2,s_i)\,R_{tu}-h_1(s_i\xi_i)\big]P-\upsilon_i\,\xi_i\,Q_{tu}.\label{eq:gamma2}
\end{equation}
Every quantity on the right is either the TA-issued secret $s_i$ or public ($ts_2,\xi_i,\upsilon_i,R_{tu},Q_{tu}$; $ID_i$ needs $id_i$, which
the TA also holds). The UAV's own random value $r_i$ enters only through $\xi_i$, and $\xi_i$ is public.

\begin{flawbox}[title={Key escrow and no non-repudiation (Table II claims F4 and F10 do not hold)}]
\begin{finding}\label{f:escrow}
The TA (or anyone who breaches it) can produce a valid request for any UAV $i$ in any session: choose $r_i$, compute
$\xi_i,ID_i,\upsilon_i$, then compute $\Gamma_i$ by~\eqref{eq:gamma2}. The request satisfies~\eqref{eq:ga3} for the same algebraic reason an honest one does.
\end{finding}
\textbf{Why.} In real certificateless cryptography~\cite{alriyami} a user's full private key combines a KGC-issued \emph{partial} key with a
\emph{secret value} the user picks and never reveals. Here the user's part ($r_i$) is effectively published, so only the TA-issued part remains.
The scheme is effectively a shared-key scheme with the TA, presented in pairing notation.

\textbf{Consequences.}
(i) \emph{Key escrow} (F4 \no): the TA can sign for any UAV. It also computes every $\sigma_i$ in GA2, so it can read the group key $\kappa_{tu}$.
(ii) \emph{Non-repudiation} (F10 \no): verification needs $\mu_i=\hat e(h_3(\cdot)P,s_iP)$, which only the holder of $s_i$ can compute. Any
credential could therefore have been produced by the TA, and a UAV can always deny it.
(iii) The pairings are not even necessary: the TA could check $\Gamma_i$ by recomputing~\eqref{eq:gamma2} with two scalar multiplications.
\end{flawbox}

% ---------------------------------------------------------------------
\subsection{Finding 3: no forward secrecy}
\begin{flawbox}[title={One leaked $s_i$ exposes every past group key (Table II claim F7 does not hold)}]
\begin{finding}\label{f:fs}
The decrypting key $\sigma_i=h_1(s_i\xi_i)P$ is a deterministic function of the long-term secret $s_i$ and the public $\xi_i$. An adversary that
records traffic and later learns $s_i$ (for example by capturing the UAV) recomputes $\sigma_i$ for every recorded session and then
$\kappa_{tu}=\Lambda(\sigma_i)\bmod\sigma_i$ from the recorded coefficients.
\end{finding}
\textbf{Why.} Forward secrecy needs an \emph{ephemeral} secret that is erased after the session, for example a Diffie--Hellman exponent. The scheme has none.
\end{flawbox}

% ---------------------------------------------------------------------
\subsection{Finding 4: the CRT group key leaks, and revocation fails}
\paragraph{(a) The construction is not well defined.}
Step GA2 defines $\sigma_i=h_1(s_i\xi_i)P$, an \emph{elliptic-curve point}. Step GD1 then multiplies points ($\omega_i=\prod_{j\ne i}\sigma_j$), inverts
modulo a point ($\varphi_i\equiv\omega_i^{-1}\bmod\sigma_i$), and Step GD6 reduces modulo a point. None of these operations is defined for points.
Under any reading that makes them work (for example $\sigma_i$ replaced by an integer derived from the point), the paper still omits three conditions
that correctness requires: the $\sigma_i$ must be pairwise coprime, $\kappa_{tu}<\min_i\sigma_i$, and~\eqref{eq:lambda} must be evaluated over the
integers, not modulo $q$.

\paragraph{(b) The polynomial reveals everyone's key.}
\begin{lemma}\label{lem:leak}
Let $C=\kappa_{tu}\sum_i\omega_i\varphi_i$. Then for every member $j$, \;$\Lambda(\sigma_j)=C$ \;and\; $\Lambda(x)-\Lambda(\sigma_j)=\prod_{i=1}^{n}(x-\sigma_i)$.
\end{lemma}
\begin{proof}
By~\eqref{eq:lambda}, $\Lambda(x)=C+\prod_i(x-\sigma_i)$ and the product vanishes at $x=\sigma_j$. Subtracting gives the second identity.
\end{proof}
So any member $j$ knows its own $\sigma_j$, computes the constant $\Lambda(\sigma_j)$, and obtains the monic integer polynomial $\prod_i(x-\sigma_i)$.
Its roots, which are the decrypting keys of \emph{all} members, can be found in polynomial time (integer root finding). Note also that the coefficients
$a_1,\dots,a_n$ are exactly the coefficients of $\prod_i(x-\sigma_i)$; only $a_0$ is masked.

\begin{flawbox}[title={A revoked UAV still obtains the new group key (verified)}]
\begin{finding}\label{f:crt}
In Step DU1 the TUAV builds $\Lambda^*(x)=\prod_{i\in S'}(x-\sigma_i)+\kappa^*_{tu}\sum_{i\in S'}\omega_i\varphi_i$ over the remaining set $S'$, and
\emph{reuses} the remaining members' old $\sigma_i$ (the paper presents this as an advantage). A member $k$ that is removed had already learned every
$\sigma_j$ through Lemma~\ref{lem:leak}. It picks any remaining $j\in S'$ and computes
\[
  \kappa^*_{tu}=\Lambda^*(\sigma_j)\bmod\sigma_j .
\]
\end{finding}
\textbf{Verification.} Our script \texttt{scripts/crt\_revocation\_attack.py} builds GD1--GD2 for $n=12$ members with 160-bit prime moduli, recovers
all twelve $\sigma_i$ from member~0's view, revokes two members, rebuilds $\Lambda^*$ by DU1, and checks the attack:
\begin{center}
\begin{tcolorbox}[colback=black!85,colframe=black!85,coltext=white,width=0.82\linewidth,arc=1mm,boxrule=0pt,fontupper=\ttfamily\small]
correctness holds for all 12 members\\
(A) insider recovered all sigma\_i: True\\
\hspace*{1.5em}revoked member's own sigma gives kappa*? False\\
(B) revoked member, using a stolen sigma, gets kappa*: True
\end{tcolorbox}
\end{center}
\textbf{Impact.} Group forward secrecy after revocation fails, and so does the claim of Theorem~5 that ``removed UAVs are locked out''. A UAV that joins
once and is then evicted keeps reading the group traffic.
\end{flawbox}

\begin{figure}[H]
\centering
\begin{tikzpicture}[font=\sffamily\small]
  \node[ent=hoc] (tu) at (0,0) {TUAV};
  \node[draw=black!40,fill=soft,rounded corners,align=center,minimum width=4.6cm] (pub) at (5,0)
      {broadcast $\{a_0,\dots,a_n\}$\\ $\Lambda(x)=C+\prod_i(x-\sigma_i)$};
  \draw[msg=hoc] (tu) -- (pub);
  \node[ent=bad,minimum width=3cm] (k) at (11,0) {revoked UAV $k$};
  \draw[msg=bad] (pub) -- node[lbl,above]{knows $\sigma_k$} (k);
  \node[act=bad,anchor=north] (s1) at (11,-1.2) {1) $C=\Lambda(\sigma_k)$\\2) roots of $\Lambda(x)-C$\\ \;\;$\Rightarrow$ every $\sigma_j$};
  \node[draw=black!40,fill=soft,rounded corners,align=center,minimum width=4.6cm] (pub2) at (5,-3.2)
      {after DU1: new $\{a^*_k\}$\\ same $\sigma_j$ for staying members};
  \node[act=bad,anchor=north] (s2) at (11,-3.0) {3) $\kappa^*=\Lambda^*(\sigma_j)\bmod\sigma_j$\\ \textbf{new key recovered}};
  \draw[msg=hoc] (tu) |- (pub2);
  \draw[msg=bad] (pub2) -- (s2);
  \draw[msg=bad,dashed] (s1) -- (s2);
\end{tikzpicture}
\caption{Revocation break (Finding~\ref{f:crt}): an insider learns every decrypting key from the public coefficients and later uses a
staying member's key to open the updated polynomial.}
\label{fig:crtattack}
\end{figure}

\paragraph{(c) The real size of the broadcast.}
The coefficients $a_k$ are integers whose length grows with $n$ (they are products of up to $n$ 160-bit moduli). Table~V of~\cite{base} counts only
the request, $104n$ bytes. We computed the exact size of $\{a_k\}$ for random 160-bit prime moduli (Table~\ref{tab:crtsize}). Because Step GD3 sends
the acknowledgement ``one by one to each vehicle'', the total traffic is $n$ times larger again: $O(n^2)$ bits per ACK and $O(n^3)$ in total.

\begin{table}[H]
\centering\small
\caption{Measured size of the CRT coefficient set (\texttt{scripts/compute\_results.py}).}
\label{tab:crtsize}
\begin{tabular}{@{}rrrr@{}}
\toprule
$n$ & per ACK & all $n$ ACKs & paper's Table~V ($104n$) \\
\midrule
 20 &   4.1 KiB &  83.0 KiB &  2.0 KiB \\
 60 &  36.0 KiB &   2.11 MiB &  6.1 KiB \\
120 & 142.7 KiB &  16.72 MiB & 12.2 KiB \\
\bottomrule
\end{tabular}
\end{table}

% ---------------------------------------------------------------------
\subsection{Finding 5: the unforgeability proof is empty}
\begin{flawbox}[title={Theorem 1 reduces to a problem that is not hard}]
\begin{finding}\label{f:proof}
The proof of Theorem~1 embeds the CDH instance $(aP,bP)=\big(s_{tu}P,\ (\sum_i\upsilon_i\xi_i)P\big)$. But $\upsilon_i$ and $\xi_i$ are sent in clear in
every request, so $b=\sum_i\upsilon_i\xi_i$ is public and $abP=b\cdot(s_{tu}P)$ can be computed directly. Solving this ``instance'' is
not hard, so the reduction proves nothing.
\end{finding}
\textbf{Further gaps.} (i) The reduction is built around the TUAV's secret $s_{tu}$, whereas forging a \emph{UAV's} credential should depend on that
UAV's secret. (ii) Theorems~2--5 are argued only in words. (iii) The batch check~\eqref{eq:ga3} uses no random small exponents~\cite{bgr}. Because
$\prod_i\psi_i=\hat e\big(\sum_i(\Gamma_i+\sigma_i),P\big)$, replacing $\Gamma_1,\Gamma_2$ by $\Gamma_1+\Delta,\ \Gamma_2-\Delta$ leaves the batch valid
although both individual credentials are now invalid. (iv) Step GA4 has no else-branch, so a single bad request rejects the whole batch and the
culprit cannot be identified.
\end{flawbox}

% ---------------------------------------------------------------------
\subsection{Finding 6: the cost accounting leaves out the expensive part}\label{sec:f6}
\begin{flawbox}[title={The reported 151\,ms excludes the TA's pairings (and the satellite delay)}]
\begin{finding}\label{f:cost}
Section VI-A states that ``the computation overhead on the TUAV side is considered'' for GA. However, Step GA2 performs, per UAV, three pairings
($\psi_i,\mu_i,\delta_i$), three scalar multiplications in $\G_1$ and one point addition:
\[ T_{\TA}=3T_P+3T_{SM\text{-}P}+T_{PA\text{-}P}+2T_{SH}=29.7414\ \text{ms per UAV.} \]
\end{finding}
In addition, $\Gamma_i$ is costed in the 160-bit group $\G$ ($T_{SM\text{-}E}$), but GA2 pairs it, so it must live in the pairing group $\G_1$.
Correcting both (constants from~\cite{base}):
\begin{center}\small
\begin{tabular}{@{}lccc@{}}
\toprule
Quantity & reported in~\cite{base} & honest accounting & ratio\\
\midrule
UAV signing (US) & 1.8203 ms & 5.3970 ms ($\Gamma_i\in\G_1$) & $3.0\times$\\
Single verification (SA) & 2.1578 ms & 31.8992 ms + RTT & $14.8\times$\\
Group auth.\ $n{=}120$ (GA) & 151.19 ms & 3720.16 ms + RTT & $24.6\times$\\
Request size & 104 B & 192 B ($\Gamma_i\in\G_1$, 128 B) & $1.8\times$\\
Satellite traffic per UAV & not reported & 192 B up + 512 B down & --\\
Key ACK, $n{=}120$ & not reported & 142.7 KiB each & --\\
\bottomrule
\end{tabular}
\end{center}
\textbf{Impact.} With honest accounting the base scheme is the \emph{slowest} of the five schemes it compares (Mei 696\,ms, Kumar 1985\,ms, Ali 333\,ms,
Xu 1668\,ms at $n{=}120$), and it is the only one that needs a wide-area round trip.
\end{flawbox}

% ---------------------------------------------------------------------
\subsection{Summary of findings}
\begin{table}[H]
\centering\small
\caption{Our findings mapped to the base paper's claims.}
\label{tab:findings}
\begin{tabularx}{\linewidth}{@{}c>{\raggedright\arraybackslash}X l l c@{}}
\toprule
\# & What fails & Where & Claim affected & Verified\\
\midrule
1 & Authentication needs a live satellite link to the TA & GA1--GA3 & ``infrastructure-less'' & by inspection\\
2 & TA can forge any UAV; verification needs $s_i$ & US2, GA2 & F4, F10 & algebraic\\
3 & Past group keys recoverable from $s_i$ & GA2, GD6 & F7 & algebraic\\
4 & Every member learns all $\sigma_i$; revoked UAV gets new key & GD2, DU1 & F3, F5, Thm 5 & script\\
5 & Theorem 1's reduction is empty; batch check unsound; no failure branch & Thm 1, GA3--4 & F1, F8 & algebraic\\
6 & TA pairings and satellite RTT omitted; $\Gamma_i$ in wrong group & Sec. VI & Tables III--V & recomputed\\
\bottomrule
\end{tabularx}
\end{table}

\begin{keybox}
All six findings share one root cause: \textbf{the TA is inside every operation}. It holds every secret, performs every check, and is reachable only over the
link the disaster degrades. Removing the TA from the field is therefore the central design goal of our framework.
\end{keybox}
